\chapter{Background}\label{cp:theory} %TODO maybe call this chapter
This chapter discusses the findings of researching existing and proposed algorithms for determining systolic and diastolic blood pressure through the oscillometric method. This report focusses on relatively simple algorithms that can be implemented in a real-time application, which is the goal of this project.

The first part of this chapter defines what blood pressure is and how it is measured manually. Afterwards 
%TODO: do I inlcude python code showing why certain algorithms do not work or do I leave that be?

A brief section at the end of this chapter explores algorithms beyond the scope of this project that have been proposed in recent literature.

\section{What is Blood Pressure?}
Blood pressure is an important bio-medical measurement and often used for  diagnostics in cardiovascular diseases. Most commonly, systolic and diastolic blood pressure are mentioned in pairs, but what do those numbers mean?

The heart is a pump that has two phases. When the ventricle is relaxed, the heart fills up and the diastolic blood pressure (DBP) is observed. When the heart contracts, it pushes the blood through the arterial system and the systolic blood pressure (SBP) is observed. A third characteristic is the difference between SBP and DBP, the pulse pressure (PP). Finally, the mean arterial pressure (MAP) is defined as the average pressure in the artery. Figure \ref{fig:BP} shows how these pressure values are connected in one heart beat. MAP is the area underneath the blue curve divided by the time of one pulse, indicated by the yellow cross-hatched area. \cite{Boron2012}\paragraph{}

\begin{figure}[]
\centering
%\includegraphics[scale=0.125]{GlaLogo.pdf}
\caption{Blood pressure characteristic values. Source: \cite{Boron2012}}
\label{fig:BP}
\end{figure}


Blood pressure can be measured with invasive and non-invasive methods. The most accurate methods are invasive, but require professional expertise because a catheter has to be injected into the blood vessel. Subsequentially, non-invasive methods are more common because they are more convenient to use. With automated methods, taking blood pressure at home has become convenient and easy. However, the inaccuracies of these measurements are often ignored or even unknown.

\section{Manual Blood Pressure Measurement}
The manual method of measuring blood pressure relies on a sphygmomanometer and listening to the heart or Korokoff sounds (named after their discoverer) with a stethoscope. This is called the \emph{auscultatory method}. \emph{Palpation} is another method where the pulse is felt at the wrist. It only allows detection of the SBP. The ausculatory method is described below because it provides recommendations that are true for all blood pressure measurements using a cuff. 

The sphygmomanometer is shown in figure \ref{fig:sphy}. It consists of a inflatable cuff that is connected to a rubber pump and a scale that indicates the pressure in the cuff in \SI{}{\mmHg}. The pump is equipped with a valve that can be opened to release air.


\SI{}{\mmHg} is a unit to descibe pressure. \SI{1}{\mmHg} is defined as the pressure generated by a column of mercury of \SI{1}{\mm} heigth. Millimeters of Mercury is not part of the International System of Units (SI), but was defined by it before 2019 as \SI{1}{\mmHg} = \SI{133.322}{\Pa} using standard gravitiy.\cite{SI2008} Since 2019, \SI{}{\mmHg} is no longer included in the SI brochure. However, it is still widely used in the medical field.
 
 
\begin{figure}[h]
\centering
%\includegraphics[scale=0.125]{GlaLogo.pdf}
\caption{Picture of a sphygmomanometer showing the cuff, inflation pump and mercury meter.}
\label{fig:sphy}
\end{figure}

Intra-arteria or direct BP measurement is considered the 'gold standard', but most studies compare the developed algorithm with the ascultatory method, which underestimated SBP and overestimates DBP. \cite{Sapinski1996} Nonetheless, the protocols to test automatic blood pressure moitors given by the British Hypertension Society (BHS) and their American counterpart, the Association for the Advancement of Medical Instrumentation (AAMI), both use the auscultatory method as a reference. \cite{Jazbinsek2010,BHS,AAMI}
 
\subsection{Auscultatory Method}
It is important that the cuff size is appropriate for the patient's arm. The rubber bladder inside the cuff should cover more than \SI{80}{\%} but less than \SI{100}{\%} of the arm's circumference.

The cuff is completely deflated and tightly fit around the upper arm of the patient. The center of the bladder should be over the brachial artery and the arm supported at heart level. The stethoscope is to be placed over the artery between the cuff and the patient's elbow and should not touch the cuff. This could influence the reading because it puts additional pressure on the artery and can interfere with hearing the sounds.\cite{NHS2019}\cite{Reeves1995}

The valve at the pump is completely closed and the cuff is inflated to a pressure of about \SI{30}{\mmHg} above the expected systolic pressure. This causes the artery to flatten and not let any blood through. Subsequently, the valve is slightly opened to slowly deflate the cuff at a rate of approximately \SI{3}{\mmHg/\second}. The artery will open and let small amounts of blood through. Systolic pressure is read when a pulse is first heard. Meanwhile, deflation continues. Diastolic pressure is read when the sounds disappear completely. Ultimately, blood can flow freely thorugh the artery. After making sure that no further sounds can be heard by deflating for at least another \SI{10}{\mmHg}, the valve is opened completely to rapildly empty the cuff.   \cite{NHS2019,Reeves1995}

There are some discrepancies in literature on how to define the point of the diastolic blood pressure. Some references define it as when the Korokoff sounds dissappear completely. \cite{NHS2019,Reeves1995} Others, predominantly older ressouces, define it as the point where the sounds are muffled.\cite{Boron2012} Here, the former is assumed, because most current literature does and some literature suggests that the auscultatory method overestimates diastolic blood pressure. \cite{Chandrasekhar2019} 

\subsection{Calculation of MAP}
Traditionally, MAP is calculated as an estimate by adding a third of the PP to the DBP. Or, how it is more commonly reffered to, adding SBP and twice the DBP and dividing the result by 3 (equation \ref{eq:MAP}). According to calculations by Bos \textit{et al.}, this generally underestimates the MAP and they suggest to add $40\%$ rather than a third of DBP. Furthermore, they suggest to use the MAP measured by oscillometric devices rather than calculating it from values obtained through the auscultatory method. \cite{Bos2007}

\begin{equation}
MAP = \frac{SBP-DBP}{3}+DBP = \frac{SBP+2 DBP}{3}
\label{eq:MAP}
\end{equation}

According to Joe \cite{Joe2019} nurses in intensive care units (ICU) today still use the manual calculation method on the automatically obtained SBP and DBP rather than rusting the measured MAP, which introduces large errors. %https://www.youtube.com/watch?v=TG0NcqKeRWE

\section{Automatic Blood Pressure Measurement}

Most modern automatic devices use the oscillometric method to detmine blood pressure. They are used both at home and in professional environments. 

Automatic blood pressure monitors measure pressure in a cuff similar to the one used in the manual method and use the small pressure changes (oscillations) extracted from the deflation curve to estimate blood pressure. Commercially available devices are usually equipped with an electic pump and the bladder contains a pressure sensor. Otherwise, they look similar to the cuff used for manual measurements. Some devices use wrist cuffs, but those are rarely recommended. \cite{BIHSOC2020} The most significant problem using wrist measurements is the influence of gravity on blood pressue and the accompaning risk of systematic errors when the measurement is not taken at heart level. \cite{Boron2012} This problem is avoided by taking blood pressure at the upper arm. 

\paragraph{General Procedure} The procedure is also similar to the manual method. The pressure in the cuff is increased to a level above the expected systolic pressure, which leads to the artery being pushed close. While releasing the pressure slowly through the valve, blood starts flowing through the artery, resulting in small increases in pressure realative to the continuous deflation of the cuff. As the cuff deflates further, these relative changes increase to a maximal value before they start decreasing agian.\cite{Forouzanfar2015,Drzewiecki1994,Ursino1996} The shape, magnitude or envelope of these oscillations are used in various automatic blood pressure algorithms. Methods that do not use oscillations exist, but are not discussed here.

\paragraph{Oscillation Extraction} There are two ways the oscillatilometric waveform (OMW) is extracted from the deflating pressure curve. The first one is filtering. A  bandpass filter with a lower cut-off frequency between \SIrange{0.1}{0.5}{\Hz} and a upper cut-off frequence between \SIrange{0.1}{0.5}{\Hz} is recommendded. \cite{Forouzanfar2015} Implementations usually use first \cite{Lim2015} to sixth order \cite{Jazbinsek2010} Butterworth filters. They are known for their flat pass-band response and a good frequency response. 

The second approach is to use detrending. It requires to know the beginning of each pulse to be able to fit a line of continuously deflating pressure to the identified points. Its advantage is, that it additionally reproduces an esimated defaltion curve. However, it requires additional data, for example ECG, for the pulse detection.\cite{Forouzanfar2015}

\paragraph{Envelope} The envelope of the OMW is used by basic automatic algorithms. Formation of the envelope is achieved in different ways. The simpleast way is to register only local maxima. Similar to that, and most common, is the subtraction of the following through from a local peak or interpolating the curve of local maxima and local minima to subtract them from each other. Some algorithms additionally fit a courve on the obtained oscillometric waveform envelope (OMVE) in an attempt to remove artifacts. \cite{Forouzanfar2014}


\subsection{Maximum Amplitude Algorithm}
The maximum amplitude algorithm (MAA) is based on the assumtion that the oscillations are maximal when the pressure in the cuff equals arterial pressure. Accordingly, the recorded pressure at which oscillations are maximal is considered a valid estimate of the MAP \cite{Babbs2012,Geddes1982,Drzewiecki1994}\cite{Ramsey1979}, as long as the compression chamber is kept small \cite{Mauck1980}. To avoid introducing errors, the cuff should always be tightly fit, because air volume in the cuff causes maximum oscillations to be above the true MAP. Hence, Ursio suggestes to use the lowest pressure of the platau of maximal oscillations as the value to assess MAP.\cite{Ursino1996}

A recent publication by Chandrasekhar \textit{et al.} \cite{Chandrasekhar2019} employs a more complex model than the one originally introduced by Mauck \textit{et al.}\cite{Mauck1980}. They conclude that the MAA results in a weighted average of systolic and disastolic BP. According to their model, MAA underestimates MAP for higher blood pressure. 

\paragraph{Fixed-Ratio Algorithm} The fixed-ratio algorithm builds on the MAA and is based on the assumption that the systolic and diasotlic blood pressure occure at specific fractions of the maximum oscillations before and after its occurence. Determing SBP and DBP after appliyng the MAA was tested by Geddes \textit{et al.} in 1982 \cite{Geddes1982}. Geddes recorded Korokoff sounds while measuring oscillations in order to find a ratio of oscillation amplitudes for SBP and DBP. They defined the ratio for the systolic pressure to be 0.5 and the ratio for the diastolic pulse as 0.8. These ratios were found in an experimental way and Geddes acknowledges that the systolic pressure is overestimated and the ratio for the diastolic pulse is not constant for a range of different diastolic pressures.

Later studies tried to find accurate ratios, mostly experimentally with the ratio for SBP usually beeing determined between 0.45 and 0.73 and the ratio for DBP a bit higher between 0.69 and 0.83.\cite{Drzewiecki1994,Forouzanfar2015} Mathematical models confirmed that a generalised ratio cannot be found. Parameters like the age of the subject and in this regard, the arterial stiffness as well as pulse pressure influence the ratios. \cite{Ursino1996} A higher PP results in a smaller ratio for SBP and larger ratio for DBP. A stiff arterial wall causes volume changes to happen slower. While this has a small effect on the systolic ratio, the diastolic ratio decreases with the stiffness of the artery.\cite{Babbs2012}

% volume changes les rapidly: artery is stiffer!babbs2012
% pulse pressure: higher pp: smaller s ratio, larger d ratio

%- algorithms are usually adjusted to match readings of the auscultatory method, but that has its own shortcommings and might introduce further errors (some use other oscillometric devices which is even worse)

\subsection{Derrivative Algorighm}
The derrivative algorithm uses the OMVE and plots it against the defation pressure. The points of maximal and minimal slope are determined. The pressure point where the derivative of the OMVE is maximal is assumed to be the diastolic BP and where the derivative is minimal is the systolic BP respectively. \cite{Jazbinsek2010,Forouzanfar2015} 

Even though mathematical models have confimed this method to produce valid estimates of both SBP and DBP without the need for empirically obtained ratios, it is extremely vulnerable to noise and therefore not used in practical applications.  \cite{Babbs2012,Chandrasekhar2019}

%Derrivative
%pulse morphology
%model based
%NN
%PTT
%
%mention artifact remvoal:
%what artifacts: high frequency noise, movement, muscle contractions
% --> how to make sure only artifacts and not information is remvoed?



\subsection{Other Algorithms}
Since the above mentioned algorithms use only the envelope of the oscillations to determine BP, 

%TODO evaluation of stiffnes with single pulse

using pulse morphology and self defined indeces, in the end they settle for using 'OPI': taking the '1st' derrivative of each pulse and taking the maximal value of that to form another slope that looks similar to the original oscillations anyway and then again using ratios to get SBP and DBP. then comparing it to measurements from an omron device. uses datasets with almost 100 pulses (mine work with 10 pulses) \cite{Mafi2012,Mafi2011}

%TODO move to discussion?
using a cuff equipped with a microphone to record Korotkoff sound to detect SBP and DBP but only to then again define a ratio to calculate blood pressure from the oscillation amplitude. Also using the derrivative algorithm, but with many constraints so that a result can be obtained. rat\cite{Jazbinsek2010,Jazbinsek2005,Jazbinsek2016}


a lot of $new$ algorithms can be brought back to the MAA method, they use NN or models to improve the ratios, but in the end, it is the same.